%% This is a JAIR Example File Compiled by Nicholas Mattei (nsmattei@tulane.edu) 
%% and Odd Erik Gundersen (odderik@ntnu.no)
%% and Mykel Kochenderfer (mykel@stanford.edu)
%% September 15, 2025
%%
%% This file is based off the ACM Latex Template https://www.acm.org/publications/proceedings-template
%% Revision 2.12 (12/28/2024)
%% 
%% Please see https://www.jair.org/index.php/jair for more information and submission instructions.
%%

%% The first command in your LaTeX source must be the \documentclass
%% command.
%%
%% For final camera ready please change the
%% command to \documentclass[]{jair}.
%%
\documentclass[manuscript, screen, review]{jair}
\usepackage{caption}
\usepackage{subcaption}
\usepackage{algorithm}
\usepackage{algorithmic}
\setcopyright{cc}
\copyrightyear{2025}
\acmDOI{10.1613/jair.1.xxxxx}

%%
\JAIRAE{Insert JAIR AE Name}
\JAIRTrack{} % Insert JAIR Track Name only if part of a special track
\acmVolume{}
\acmArticle{}
\acmMonth{12}
\acmYear{2025}

\RequirePackage[
  datamodel=acmdatamodel,
  style=acmauthoryear,
  backend=biber,
  giveninits=true,
  uniquename=init
  ]{biblatex}

\renewcommand*{\bibopenbracket}{(}
\renewcommand*{\bibclosebracket}{)}

%%
%% For managing citations, use BibLaTeX with the acmauthoryear style.
%% The next line specifies the bibliography file.
\addbibresource{sample-base.bib}

%%
%% end of the preamble, start of the body of the document source.
\begin{document}

%%
%% The "title" command has an optional parameter,
%% allowing the author to define a "short title" to be used in page headers.
%\title{JAIR Example Template}
\title[NELE]{NELE: Non-Uniform Rational B-Spline Elementwise Learnable}

%%
%% The "author" command and its associated commands are used to define
%% the authors and their affiliations.
%% Of note is the shared affiliation of the first two authors, and the
%% "authornote" and "authornotemark" commands
%% used to denote shared contribution to the research and/or corresponding author.
\author{Vishakh Begari}
\orcid{0009-0008-9399-8581}
\email{vibedjentle@gmail.com}
\affiliation{%
  \institution{Independent Researcher}
  \city{Gauting}
  \state{Bavaria}
  \country{Germany}
}

%% The short list of authors must be made of the list of all authors' lastnames.
\renewcommand{\shortauthors}{Begari}
%% If this is too long and overlaps other information printed in the page headers, use
%\renewcommand{\shortauthors}{Xu et al.}

%%
%% The abstract is a short summary of the work to be presented in the
%% article.
\begin{abstract}
{\bf Background:} 
    The design of activation functions (AF) 
    is central to the expressiveness and efficiency of neural networks. 
    Yet the pursuit of a universal activation function, one that can 
    automatically adapt its shape to suit a wide variety of tasks, remains 
    an open challenge.  
    
    {\bf Objectives:}
    In this work, the author proposes \texttt{NELE} (NURBS Elementwise Learnable), 
    a learnable activation function based on Non-Uniform Rational B-Splines (NURBS).

    {\bf Methods:}
    Unlike conventional fixed activations, \texttt{NELE} 
    learns a smooth, data driven mapping from inputs to outputs 
    through a set of control points and corresponding weights.
    This yields a continuous, differentiable family of functions capable 
    of modeling a broad spectrum of nonlinear behaviors, effectively 
    unifying diverse activation patterns within a single formulation.

    {\bf Results:}
    \texttt{NELE} was evaluated on several state-of-the-art benchmark datasets.
    The empirical results indicate that \texttt{NELE} achieves consistently better performance 
    than multiple existing AF approaches.

    {\bf Conclusions:}
    The theoretical universality of NURBS, its ability to approximate a wide class 
    of functions and its inherent flexibility, positions \texttt{NELE} as a promising candidate 
    for a universal activation function.
\end{abstract}

%\begin{abstract}
%      A clear and well-documented \LaTeX\ document is presented as an
%  article formatted for publication by ACM in a conference proceedings
%  or journal publication. Based on the ``acmart'' document class, this
%  article presents and explains many of the common variations, as well
%  as many of the formatting elements an author may use in the
%  preparation of the documentation of their work.
%\end{abstract}


%% JAIR Note: 
%% Do not include ACM CCS Concepts or Keywords


%% To be updated by authors.
\received{26 December 2025}
\received[accepted]{26 December 2025}

%%
%% This command processes the author and affiliation and title
%% information and builds the first part of the formatted document.
\maketitle

\section{Introduction}
Activation functions (AF) are a core component of artificial neural networks, determining how neurons process input 
signals and introduce non-linearity into the network. Their evolution reflects the growing complexity 
and depth of neural models: from the early binary step functions of McCulloch-Pitts \cite{mcculloch1943logical}
neurons and perceptrons, 
which could handle only simple, linearly separable tasks, to smooth, differentiable functions like \texttt{sigmoid} \cite{MINAI1993845}
and \texttt{tanh} \cite{6795724} that enabled backpropagation in multilayer networks. The advent of Rectified Linear Unit (\texttt{ReLU}) \cite{relu} and its variants solved 
vanishing gradient issues and facilitated deep learning, while modern activations like Exponential Linear Unit (\texttt{ELU}) \cite{clevert2016fastaccuratedeepnetwork}, 
\texttt{Swish} \cite{ramachandran2017searchingactivationfunctions}, 
\texttt{Mish} \cite{misra2020mishselfregularizednonmonotonic} 
further enhance network performance by improving training stability and modeling complex patterns.
While classical activations such as \texttt{ReLU}, \texttt{tanh}, and Gaussian Error Linear Unit (\texttt{GELU}) \cite{hendrycks2023gaussianerrorlinearunits}
have enabled remarkable success, each represents only a narrow subset of possible nonlinear behaviors. 
Polynomial activation functions are unbounded and prone to numerical instability, 
often requiring high degrees or large networks to achieve sufficient expressiveness \cite{Hornik1989Multilayer}
\cite{Leshno1993NonPolynomial} \cite{Goodfellow2016DeepLearning}.

The concept of spline-based AF is not new. Early studies, such as those by 
\cite{guarnieri1999multilayer}, introduced adaptive spline activations to provide flexible, 
data-driven nonlinearities within neural networks. 
More recent works, including \cite{hryniowski2018polyneuronautomaticneurondiscovery} and 
\cite{bohra2020learning}, extended this idea by learning spline parameters directly during training, 
allowing the AF itself to adapt to the data. 
Building on this foundation, Non-Uniform Rational B-Splines (NURBS) can be seen as a natural 
extension, offering even greater flexibility through rational weighting and non-uniform knot placement. 
This enables NURBS based activations to represent a broader class of nonlinear mappings while 
retaining the smoothness and local control properties of traditional splines.
Using appropriate parameters, existing activation functions can be approximated using the NURBS based approach.

\section{Theory}
NURBS \cite{nurbs} are a mathematical representation used to describe curves and surfaces 
in a flexible and precise way. They generalize B-splines and Bézier \cite{Prautzsch2002Bezier} curves 
by introducing non-uniform knot vectors and rational weighting, 
allowing the representation of both free-form shapes and 
exact analytic geometry. The non-uniform spacing of knots enables localized shape control, 
while rational weights allow exact representation of conic sections 
and other geometries not achievable with polynomial B-splines.
A NURBS curve of degree $p$ is defined as:
\begin{equation}\label{equation:nurbs}
\mathbf{C}(u)=
\frac{
\displaystyle\sum_{i=0}^{n} N_{i,p}(u)\, w_i\, \mathbf{P}_i
}{
\displaystyle\sum_{i=0}^{n} N_{i,p}(u)\, w_i
},
\qquad u \in [u_0, u_m].
\end{equation}
where, $P_i$ = control points, $w_i$ = weights associated with the control points, 
$N_{i,p}(u)$ = B-spline basis functions of degree $p$, ${u_0, u_1,...u_m}$ = knot vector, 
$n + 1$ = number of control points and $m = n + p + 1$ = knot count relationship.

Setting all the weights equal ($w_i$ = 1), Equation~\ref{equation:nurbs} reduces to 
$C(u) = \displaystyle\sum_{i=0}^{n} N_{i,p}(u)P_i$, which is a B-spline curve.
Stone-Weierstrass theorem \cite{stone1948generalized} establishes that any continuous real valued 
function on a compact interval can be uniformly approximated arbitrarily 
closely by elements of a suitably rich algebra of functions 
(for example, polynomials), spline spaces being piecewise polynomial 
subalgebras, inherit the same dense approximation property 
in the space of continuous functions. Since NURBS curves 
generalize B-splines by adding weights 
(and thereby include B-spline curves as a special case), the ability to represent exactly all 
polynomials up to a given degree and retain the capacity to approximate any 
continuous (and sufficiently smooth) function arbitrarily well.

For a smooth curve $y = f(x)$, the curvature is given by:
\begin{equation}
\kappa = \frac{\left|y''\right|}{\left(1 + (y')^2\right)^{3/2}}
\end{equation}
The maximum curvature occurs at points where $\frac{d\kappa}{dx} = 0$.
The solutions for $x$ corresponding to $\kappa' = 0$ for different AF are presented
in Table ~\ref{tab:AF}.
Since curvature is a geometric invariant
\cite{doCarmo1976}, let $A(u) = \displaystyle\sum_{i=0}^{n} N_{i,p}(u)\, w_i\, \mathbf{P}_i$, \quad 
$W(u) = \displaystyle\sum_{i=0}^{n} N_{i,p}(u)\, w_i$, then
\begin{equation}
\mathbf{C'}=\frac{AW' - A'W}{W^2}
\end{equation}
\begin{equation}
\mathbf{C''}=\frac{W(A''W - AW'') - 2W'(AW' - A'W)}{W^3}
\end{equation}
\begin{equation}\label{eq:nurbscurv}
  \kappa = \frac{W^3(W(A''W - AW'') - 2W'(AW' - A'W))}{(W^4 + (AW' - A'W)^2)^\frac{3}{2}}
\end{equation}
The upper bound for Equation~\ref{eq:nurbscurv} is obtained by setting the denominator to zero.
\begin{equation}
(W^4 + (AW' - A'W)^2)^\frac{3}{2} = 0
\end{equation}
\begin{equation}
  W^4 + (AW' - A'W)^2 = 0
\end{equation}
Sum of two non-negative numbers is zero only if each term is zero individually,
$W^4 = 0$, $(AW' - A'W)^2 = 0$. Simplifying gives, $W = 0$, and $A'W = AW'$. 
Substituting, $W = 0$, $AW' = 0$. In a generic cubic Bézier curve with arbitrary weights and distinct control points, 
$A \neq 0$, in general. In that case, for $n = 3$, $W'$ is linear. This implies that the upper bound of curvature, $\kappa$,
is directly proportional to the weight of the middle control point. Furthermore, if the 
x coordinates of the 3 control points are fixed, 
this results in a total of four unknown parameters that can be learned.
Figure ~\ref{fig:nurbs} illustrates how varying the weight of the middle control 
  point affects the curve: a lower weight ($w=0.2$) pulls the curve away 
  from the point, uniform weight ($w=1.0$) yields the standard quadratic B-spline, 
  and a higher weight ($w=5.0$) attracts the curve toward the control point. 
  Red markers indicate the control points used for all curves. 
\begin{figure}[H]
  \begin{center}
    \includegraphics[width=0.4\linewidth]{../../PICS/nurbs.pdf}
  \caption{NURBS curves generated from three control points with quadratic 
  degree ($p=2$) and different weights.}
  \Description{Different NURBS curves.}
  \label{fig:nurbs}
  \end{center}
\end{figure}
\begin{table}
  \caption{AF with their peak curvature values and the input positions at which these maxima occur.}
  \label{tab:AF}
  \begin{tabular}{@{}lccc@{}}
    \toprule
    AF & $y$ & $x_{\kappa' = 0}$ & $\kappa_{\text{max}}$\\
    \midrule
    \texttt{Tanh}~\ref{app:tanh} & $\frac{e^{x} - e^{-x}}{e^{x} - e^{-x}}$ & 0.9196 & 0.5073\\
    \texttt{Sigmoid}~\ref{app:sigmoid} & $\frac{1}{1 + e^{-x}}$ & 1.3819 & 0.0924\\
    \texttt{Softplus}~\ref{app:softplus} & $\ln{({1 + e^{x}})}$ & -0.5 & 0.1913\\
    \texttt{ELU}~\ref{app:elu} & $
    \begin{cases}
      x & \text{if } x \geq 0, \\
      \alpha (e^x - 1) & \text{if } x < 0
    \end{cases}$ & $-\frac{1}{2} \ln(2 \alpha^2)$ & $\frac{2}{3\sqrt3}$\\
    \texttt{SiLU}~\ref{app:silu} & $x \cdot \frac{1}{1 + e^{-x}}$ & -0.3768 & 0.4037 \\
    \texttt{GELU}~\ref{app:gelu} & $x \phi(x)$ & -0.4773 & 1.3579\\
    \texttt{ReLU} & $
    \begin{cases}
    x, & \text{if } x > 0, \\
    0, & \text{if } x \le 0.
    \end{cases}$ & undefined & undefined\\
    \texttt{Leaky-ReLU} & $
    \begin{cases}
    x, & \text{if } x > 0, \\
    \alpha x, & \text{if } x \le 0.
    \end{cases}$ & undefined & undefined\\
    \texttt{LeLU}~\ref{app:LeLU} & $
    \begin{cases}
    x, & \text{if } x > 0, \\
    e^{(1 - \beta)x} - 1 + \beta x, & \text{if } x \le 0.
    \end{cases}$ & $- \frac{\sqrt{2}}{2}\ln2$ & 0.3924\\
    \texttt{Mish}~\ref{app:mish} & $x \cdot \tanh\!\left(\ln(1 + e^{x})\right)$ & -0.4271 & 0.5900\\
  \bottomrule
\end{tabular}
\end{table}
\section{Formulation}
The motivation for \texttt{NELE} was 
\begin{itemize}
\item Universality : To enable an AF which opens up a search 
space for flexibility.
\item Attenuation : Prevent large positive values from weakening \cite{elfwing2017sigmoidweightedlinearunitsneural}.
\item Non Zero-Gradient : Produce negative side activations.
\item Smoothnesss : Maintain continuity and differentiability throughout.
\end{itemize} 

The universality follows naturally from its 
NURBS-based approximation shown in Appendix~\ref{app:nurbsparams}. Attenuation prevention with non-zero gradients 
is achieved through a piecewise formulation in which the 
positive input region preserves the identity mapping of \texttt{ReLU}, 
while the negative region is shaped to emulate smooth activations such as 
\texttt{SiLU} or \texttt{GELU}. To ensure continuity and differentiability 
across the junction of these regions, the control points and their 
associated weights in the vicinity of the origin are appropriately adjusted, 
enforcing smooth transitions between the linear and nonlinear segments.

A cubic spline is employed as it provides sufficient expressivity while maintaining low 
computational complexity and numerical stability.
Four control points ($p_0$, $p_1$, $p_2$, $p_3$) are selected, corresponding to a cubic spline, 
which represents the lowest order capable of modeling nonlinear curvature while avoiding 
overfitting and their associated weights ($w_n$). The implementation is summarized in Algorithm ~\ref{alg:nele_training}.
Continuity condition forces $p_3 = [0.0, 0.0]$ and linearity of the positive side 
sets the slope between $p_2$ and $p_3$ = 1 to maintain tangency as depicted in Figure ~\ref{fig:nurbsregen}.
A mask is added which returns 0 for $x < p_0$.
The remaining control points and their associated weights were either treated as 
learnable parameters or fixed. 
For the fixed variant, setting $p_0$ = [-4.0, 0.0], $p_1 = [0.1, 0.1]$ and $w_n = 1$ yielded the best results. 
\begin{figure}[H]
  \begin{center}
    \includegraphics[width=0.5\linewidth]{../../PICS/Approx_gelu.pdf}
  \caption{GELU approximated for $x \leq 0$ using cubic NURBS with 4 control points and optimized weights.
  NURBS curve used in experiments is also shown for comparison.}
  \Description{GELU activation function recreation.}
  \label{fig:nurbsregen}
  \end{center}
\end{figure}

\begin{algorithm}
\caption{NELE activation function}
\label{alg:nele_training}
\begin{algorithmic}[1]
\REQUIRE Training data $\mathcal{D}$, initial parameters $p_0$, $p_1$, $p_2$, $p_3$, $w_0$, $w_1$, $w_2$, $w_3$
\ENSURE Optimized parameters $\theta^*$
\IF{Learnable parameters}
    \STATE Replace $x_2$ and $y_2$ with $l/ \sqrt(2)$
    \STATE Initialize $l$, $p_0$, $p_1$, $p_3$
    \STATE Set $w_0$, $w_3$
\ELSE
    \STATE Set $p_0$, $p_1$, $p_2$, $p_3$, $w_0$, $w_1$, $w_2$, $w_3$
\ENDIF
\FOR{each batch}
    \IF{Clamped}
      \STATE Set $x_0$ = $min(\mathcal{D})$
    \ENDIF
    \STATE Evaluate $N_0 = (1-t)^3$, $N_1 = 3t(1-t)^2$, $N_2 = 3t^2(1-t)$, $N_3=t^3$, where $t \in [0,1]$
    \STATE Calculate $\text{Numerator\_y} = \sum N_n * w_n * y_n$, $\text{Numerator\_x} = \sum N_n * w_n * x_n$, $\text{Denominator} = \sum N_n * w_n$
    \STATE Calculate ${curve\_y} = \text{Numerator\_y}/\text{Denominator}$. ${curve\_x} = \text{Numerator\_x}/\text{Denominator}$.
    \STATE Find ${y}$ from ${curve\_y}$ for ${curve\_x}$ = $\mathcal{D}$ values via Linear Interpolation.
    \IF{Masking}
      \IF{Double}
        \RETURN $y$ for $\mathcal{D} \leq 0$, $\mathcal{D} \geq min(\mathcal{D})$ and $\mathcal{D}$ otherwise
      \ELSE
        \RETURN $y$ for $\mathcal{D} \leq 0$ and $\mathcal{D}$ otherwise
      \ENDIF
    \ELSE
      \RETURN $y$
    \ENDIF
\ENDFOR
\end{algorithmic}
\end{algorithm}

\section{Experiments with \texttt{NELE}}
Following the construction process outlined in Figure~\ref{fig:nurbsregen}, 
\texttt{NELE} is benchmarked against a representative set 
of standard AF, \texttt{tanh}, \texttt{sigmoid}, \texttt{softplus}, 
\texttt{ELU}, \texttt{SiLU}, \texttt{GELU}, \texttt{ReLU}, 
\texttt{Leaky ReLU}, \texttt{LeLU} \cite{BIGARELLA2026122739} and 
\texttt{Mish}. Experiments are conducted on a combination of synthetic and standard benchmark datasets: 
(i)  Non-linear synthetic data (NLSD), (ii) MNIST \cite{lecun1998mnist}, (ii) MNIST autoencoder, (iv) CIFAR-10 \cite{hope2017learning}
 and (v) CIFAR-100 \cite{krizhevsky2009cifar}. 
All empirical evaluations are performed using feed-forward neural 
network architectures. The behavior of \texttt{NELE} in 
large-scale architectures, such as transformers or 
recurrent networks, is not investigated in this work.
All models are trained for a fixed number of epochs to ensure a 
fair comparison across AF under identical optimization conditions.

\subsection{Non-linear synthetic data (NLSD)}
NLSD provides a controlled environment for assessing functional expressivity.
The neural network used for curve fitting consists of a fully connected feed-forward 
architecture with a single hidden layer of 64 neurons. 
The hidden layer employs the various AF, while the output layer is linear. 
This design isolates the contribution of the activation function and 
ensures that observed differences in performance reflect the expressive 
capacity of each activation.
In schematic form, the network can be represented as:
\begin{displaymath}
x \;\xrightarrow{\text{Linear } (\text{input\_dim} \to 64)}\;
\text{AF} \;\xrightarrow{\text{Linear } (64 \to 1)}\; y
\end{displaymath}
The model was trained with Mean-Squared-Error (MSE) loss \cite{BickelDoksum2015MathStats}.
MSE provides a smooth, differentiable penalty for prediction errors, 
emphasizes larger deviations, and produces outputs that closely follow 
the target curve \cite{Bishop2006PRML}, which is important for subsequent smoothness analysis.
To evaluate the performance of different AF for function approximation, 
a series of synthetic datasets representing diverse nonlinear relationships is generated. 
Specifically, the target outputs $y$ were constructed as follows:
\begin{equation}
y_1 = \sin(x) + \epsilon\tag{a}
\end{equation}
\begin{equation}
y_2 = \sin(x) + 0.5 \cdot \sin(3x) + \epsilon\tag{b}
\end{equation}
\begin{equation}
y_3 = e^{-0.5 x} + \epsilon\tag{c}
\end{equation}
\begin{equation}
y_4 = \tanh(x) + \epsilon\tag{d}
\end{equation}
\begin{equation}
y_5 = x^2 +  \epsilon\tag{e}
\end{equation}
\begin{equation}
y_6 = \frac{x^3}{e^{x} - 1 + 10^{-6}} + \epsilon\tag{f}
\end{equation}
where $\epsilon \sim \mathcal{N}(0,1)$ denotes Gaussian noise added to simulate measurement 
or environmental variability. The chosen functions span a range of behaviors, 
including oscillatory patterns, exponential decay, saturating nonlinearities 
and polynomial growth, thereby providing a comprehensive testbed for assessing the 
generalization capability of different neural network architectures. 
The networks were evaluated based on their ability to fit the underlying 
function while mitigating overfitting to noise, $\epsilon = 0.2$, chosen relative to the scale of the input data. Each model was trained for 
a maximum of 300 epochs with a learning rate ($10^{-2}$, $10^{-3}$, $10^{-4}$) and parameter initialisation 
optimised to give the lowest training loss through grid search.
The studies were repeated across 7 runs with different random seeds to 
account for stochasticity in network initialization and data shuffling \cite{Dushatskiy_2022}.
and the median training loss across these runs is shown in 
Figure~\ref{fig:Lossplottrend}.
Preliminary results indicate that AF with smooth, bounded characteristics 
(e.g., \texttt{Tanh}, \texttt{Sigmoid}) excel in capturing oscillatory 
and saturating patterns.
These findings highlight the importance of matching AF properties 
to the underlying target function characteristics in curve-fitting tasks.
Given a set of discrete points $\{(x_i, y_i)\}_{i=1}^N$ representing a curve, 
smoothness was quantified using curvature based criteria:
For a parametric curve $(x(t), y(t))$, the curvature at a point is defined as
\begin{equation}
\kappa(t) = \frac{|x'(t)y''(t) - y'(t)x''(t)|}{\left( x'(t)^2 + y'(t)^2 \right)^{3/2}}
\end{equation}

The total curvature is approximated for discrete points using finite differences:
\begin{equation}
K_\text{total} = \sum_{i=1}^{N} \kappa_i
\end{equation}
\begin{figure}[H]
    \centering
    % Row 1
    \begin{subfigure}{0.32\textwidth}
        \centering
        \includegraphics[width=\linewidth]{../../PICS/NLSD/exp/curve_fitting.pdf}
        \caption{Exponential + Noise}
    \end{subfigure}
    \hfill
    \begin{subfigure}{0.32\textwidth}
        \centering
        \includegraphics[width=\linewidth]{../../PICS/NLSD/hyp/curve_fitting.pdf}
        \caption{Hyperbola + Noise}
    \end{subfigure}
    \hfill
    \begin{subfigure}{0.32\textwidth}
        \centering
        \includegraphics[width=\linewidth]{../../PICS/NLSD/quad/curve_fitting.pdf}
        \caption{Quadratic + Noise}
    \end{subfigure}

    % Row 2
    \begin{subfigure}{0.32\textwidth}
        \centering
        \includegraphics[width=\linewidth]{../../PICS/NLSD/sine/curve_fitting.pdf}
        \caption{Sinusoidal + Noise}
    \end{subfigure}
    \hfill
    \begin{subfigure}{0.32\textwidth}
        \centering
        \includegraphics[width=\linewidth]{../../PICS/NLSD/trig/curve_fitting.pdf}
        \caption{Trigonometric + Noise}
    \end{subfigure}
    \hfill
    \begin{subfigure}{0.32\textwidth}
        \centering
        \includegraphics[width=\linewidth]{../../PICS/NLSD/exppoly/curve_fitting.pdf}
        \caption{Polynomial + Noise}
    \end{subfigure}

    \caption{Curve fitting results for different target functions 
    using the proposed neural network, illustrating the network's 
    ability to approximate diverse nonlinear patterns.}
    \label{fig:curvefittingtrend}
    \Description{Non linear synthetic data curve fitting}
\end{figure}

\begin{table}
  \caption{Comparison of smoothness metrics for different activation functions and \texttt{NELE} tuned for lowest training loss. 
The table reports total curvature ($K_\text{total}$) for each AF. 
Lower values indicate smoother behavior.}
  \label{tab:smooth}
  \begin{tabular}{@{}llllllll@{}}
    \toprule
    AF & (a) & (b) & (c) & (d) & (e) & (f) \\
    \midrule
    \texttt{Tanh} & 46.72 & 39.51 & 41.47 & 73.0 & 155.65 & 35.13 \\ 
    \texttt{Sigmoid} & 26.86 & 31.5 & 44.09 & 63.14 & 68.48 & 27.32 \\ 
    \texttt{Softplus} & 31.96 & 34.94 & 39.45 & 55.88 & 57.04 & 23.04 \\ 
    \texttt{ELU} & 45.33 & 33.36 & 39.77 & 57.36 & 61.63 & 23.44 \\ 
    \texttt{SiLU} & 46.02 & 32.34 & 39.51 & 57.38 & 64.14 & 30.44 \\ 
    \texttt{GELU} & 77.24 & 34.76 & 41.84 & 61.93 & 107.0 & 50.47 \\ 
    \texttt{ReLU} & 107.69 & 55.85 & 51.58 & 75.75 & 79.89 & 51.41 \\ 
    \texttt{Leaky-ReLU} & 96.3 & 52.63 & 47.37 & 66.68 & 79.93 & 34.55 \\ 
    \texttt{LeLU} & 27.58 & 31.24 & 39.82 & 69.49 & 111.18 & 19.86 \\ 
    \texttt{Mish} & 60.13 & 34.1 & 43.28 & 59.05 & 60.72 & 28.71 \\ 
    \texttt{NELE} & 44.02 & 38.33 & 41.4 & 61.62 & 114.7 & 30.81 \\ 
  \bottomrule
\end{tabular}
\end{table}

Among the AF evaluated, the proposed \texttt{NELE} activation 
consistently demonstrates comparable and sometimes superior performance across all benchmark functions. 
Quantitatively, networks employing \texttt{NELE} achieve lower training loss and smaller 
standard deviation of curve-fitting errors relative to alternative activations, 
indicating both improved accuracy and enhanced reliability.

Following the recommended procedure for \texttt{LELU}, the neural network should 
be trained via backpropagation to minimize the mean absolute error (MAE) 
between the predicted values and the true function at the training points, 
followed by computation of diffusion terms and application of the diffusion loss. 
In the experiments, however, only the standard MSE loss was used. 
Under these conditions, \texttt{NELE} produced  
predictions with smoothness comparable to \texttt{LeLU} and training 
loss values that were lower or similar across all benchmark functions.
\begin{figure}[H]
    \centering

    % Row 1
    \begin{subfigure}{0.32\textwidth}
        \centering
        \includegraphics[width=\linewidth]{../../PICS/NLSD/exp/Loss_Bar_Chart_with_Error_Bars.pdf}
        \caption{Exponential + Noise}
    \end{subfigure}
    \hfill
    \begin{subfigure}{0.32\textwidth}
        \centering
        \includegraphics[width=\linewidth]{../../PICS/NLSD/hyp/Loss_Bar_Chart_with_Error_Bars.pdf}
        \caption{Hyperbola + Noise}
    \end{subfigure}
    \hfill
    \begin{subfigure}{0.32\textwidth}
        \centering
        \includegraphics[width=\linewidth]{../../PICS/NLSD/quad/Loss_Bar_Chart_with_Error_Bars.pdf}
        \caption{Quadratic + Noise}
    \end{subfigure}

    % Row 2
    \begin{subfigure}{0.32\textwidth}
        \centering
        \includegraphics[width=\linewidth]{../../PICS/NLSD/sine/Loss_Bar_Chart_with_Error_Bars.pdf}
        \caption{Sinusoidal + Noise}
    \end{subfigure}
    \hfill
    \begin{subfigure}{0.32\textwidth}
        \centering
        \includegraphics[width=\linewidth]{../../PICS/NLSD/trig/Loss_Bar_Chart_with_Error_Bars.pdf}
        \caption{Trigonometric + Noise}
    \end{subfigure}
    \hfill
    \begin{subfigure}{0.32\textwidth}
        \centering
        \includegraphics[width=\linewidth]{../../PICS/NLSD/exppoly/Loss_Bar_Chart_with_Error_Bars.pdf}
        \caption{Polynomial + Noise}
    \end{subfigure}

    \caption{Standard deviation of prediction errors for different target functions, illustrating 
    the variability in the neural network's curve fitting performance. Despite \texttt{LeLU} 
    reaching a lower training loss for the sinusoidal function, 
    the corresponding fit lacks smoothness.}
    \label{fig:standarddevtrend}
    \Description{Predictions}
\end{figure}
\begin{figure}[H]
    \centering

    % Row 1
    \begin{subfigure}{0.32\textwidth}
        \centering
        \includegraphics[width=\linewidth]{../../PICS/NLSD/exp/training_loss.pdf}
        \caption{Exponential + Noise}
    \end{subfigure}
    \hfill
    \begin{subfigure}{0.32\textwidth}
        \centering
        \includegraphics[width=\linewidth]{../../PICS/NLSD/hyp/training_loss.pdf}
        \caption{Hyperbola + Noise}
    \end{subfigure}
    \hfill
    \begin{subfigure}{0.32\textwidth}
        \centering
        \includegraphics[width=\linewidth]{../../PICS/NLSD/quad/training_loss.pdf}
        \caption{Quadratic + Noise}
    \end{subfigure}

    % Row 2
    \begin{subfigure}{0.32\textwidth}
        \centering
        \includegraphics[width=\linewidth]{../../PICS/NLSD/sine/training_loss.pdf}
        \caption{Sinusoidal + Noise}
    \end{subfigure}
    \hfill
    \begin{subfigure}{0.32\textwidth}
        \centering
        \includegraphics[width=\linewidth]{../../PICS/NLSD/trig/training_loss.pdf}
        \caption{Trigonometric + Noise}
    \end{subfigure}
    \hfill
    \begin{subfigure}{0.32\textwidth}
        \centering
        \includegraphics[width=\linewidth]{../../PICS/NLSD/exppoly/training_loss.pdf}
        \caption{Polynomial + Noise}
    \end{subfigure}

    \caption{Training loss curves for the neural network across different target functions, 
    showing the convergence behavior and learning dynamics for each 
    functional form and AF.}
    \label{fig:Lossplottrend}
    \Description{Training loss}
\end{figure}

\subsection{MNIST}
The performance of various activation functions was evaluated by 
replicating the experimental setup of \cite{hendrycks2023gaussianerrorlinearunits}.
Fully connected neural networks with eight hidden layers of 128 neurons each were trained
using AF listed in Table~\ref{tab:smooth}. 
All networks were trained using the Adam optimizer with cross-entropy loss 
\cite{mao2023crossentropylossfunctionstheoretical} and
learning rates of $10^{-3}$, $10^{-4}$, and $10^{-5}$. Validation used $5,000$ samples 
to select the best model. In total, seven runs were performed for 75 epochs on the complete 
dataset with the optimized hyperparameters and their median values are used for comparison.
\texttt{NELE} was stopped early at epoch 50 as it showed comparable convergence to other AF.
Figure~\ref{fig:mnistloss} shows how each AF influences convergence speed and final loss.
\begin{figure}[H]
  \begin{center}
    \includegraphics[width=1.0\linewidth]{../../PICS/MNIST/training_loss.pdf}
  \caption{Training and test accuracy trends for MNIST over epochs using different AF. Differences in the test accuracy curves reflect 
  the varying generalization performance of each AF.\label{fig:mnistloss}}
  \end{center}
  \Description{MNIST training loss}
\end{figure}
\begin{figure}[H]
  \begin{center}
    \includegraphics[width=0.5\linewidth]{../../PICS/MNIST/noise_study.pdf}
  \caption{Robustness results for MNIST over noise strengths using different AF.\label{fig:mnistnoise}}
  \end{center}
  \Description{MNIST results}
\end{figure}
\texttt{NELE} converged within 30 epochs with a 
learning rate (\textit{lr}) of $0.001$ and achieved the lowest test error~\ref{tab:mnisttesult}. 
The trained model was also tested for robustness against
inputs containing uniform noise Unif[-a,a] added to each example at different levels a, where
a = 3 is the greatest noise strength. 
Here, \texttt{NELE} shows good robustness with its performace matching \texttt{GELU} and exceeding
other AF.
\begin{table}
  \caption{Comparison of training loss and test error on the MNIST classification task
trained with different AF. \texttt{NELE} achieves the 
lowest test error among all activations evaluated.}
  \label{tab:mnisttesult}
  \begin{tabular}{@{}llcc@{}}
    \toprule
    AF &Train & Test (Noise=0) & Test (Noise=3)\\
    \midrule
    \texttt{Tanh} & 0.0076 & 2.7 & 22.94\\
    \texttt{Sigmoid} & 0.0072 & 4.45 & 16.47\\
    \texttt{Softplus}& 0.004 & 2.52 & 22.37\\
    \texttt{ELU}& 0.0056 & 2.44 & 21.51\\
    \texttt{SiLU} & 0.0045 & 2.28 & 24.64\\
    \texttt{GELU} & 0.0049 & 2.32 & 24.89\\
    \texttt{ReLU} & 0.0054 & 2.32 & 23.02\\
    \texttt{Leaky-ReLU}& 0.0044 & 2.35 & 21.22\\
    \texttt{LeLU} & 0.0002 & 2.12 & 18.5\\
    \texttt{Mish} & 0.0056 & 2.32 & 24.31\\
    \texttt{NELE}& 0.0056 & \textbf{1.83} & \textbf{25.49}\\
  \bottomrule
\end{tabular}
\end{table}

\subsection{MNIST autoencoder}
The neural network is a deep autoencoder from \cite{desjardins2015naturalneuralnetworks}
with Adam optimizer and a batch size of 64. A learning rate of $10^{-3}$ with Mean Squared Error loss 
produced the highest validation score with 5,000 samples. Figure~\ref{fig:mnistenc} shows that
with settings of \texttt{NELE} from Figure~\ref{fig:nurbsregen}, the AF
can perform quite well in tests.
\begin{figure}[H]
  \begin{center}
    \includegraphics[width=1.0\linewidth]{../../PICS/MNIST_ENC/training_loss.pdf}
  \caption{Training loss and test reconstruction error 
  trends for MNIST autoencoder over epochs using 
  different AF.\label{fig:mnistenc}}
  \Description{MNIST autoencoder results}
  \end{center}
\end{figure}

\subsection{CIFAR-10}
The neural network is a 9-layer network with the architecture and training
procedure from \cite{hendrycks2023gaussianerrorlinearunits}. The tuning was done 
using learning rates of $10^{-3}$, $10^{-4}$, and $10^{-5}$ for 100 epochs.
The best model was selected based on validation accuracy using 90-10 training/validaition split.
Seven runs were performed and the median values are reported.
Figure~\ref{fig:cifar10loss} shows the training loss and test accuracy calculated every 5 epochs for CIFAR-10 over epochs using different AF.

\begin{figure}[H]
  \begin{center}
    \includegraphics[width=1.0\linewidth]{../../PICS/CIFAR10/training_loss.pdf}
  \caption{Training and test accuracy trends for CIFAR10 over epochs using different AF.}
  \label{fig:cifar10loss}
  \end{center}
  \Description{CIFAR10 results}
\end{figure}

\subsection{CIFAR-100}
The model for CIFAR-100 is a wide residual neural network described in 
\cite{zagoruyko2017wideresidualnetworks} with learning rate scheduler and dropout keep
probability from
\cite{loshchilov2017sgdrstochasticgradientdescent} and batch normalization 
at the end of a residual block to stabilize exploding ELU's \cite{Shah_2016}
as implemented in \cite{hendrycks2023gaussianerrorlinearunits}. Setting from Figure~\ref{fig:nurbsregen} 
were used to obtain Median convergence
curves over seven runs as shown in Figure~\ref{fig:cifar100loss}.
\begin{figure}[H]
  \begin{center}
    \includegraphics[width=1.0\linewidth]{../../PICS/CIFAR100/training_loss.pdf}
  \caption{Training and test accuracy trends for CIFAR100 over 50 epochs using different AF.}
  \label{fig:cifar100loss}
  \end{center}
  \Description{CIFAR100 results}
\end{figure}
\section{Conclusions}
\texttt{NELE} exhibits a combination of properties that make it particularly 
well suited for functional approximation tasks. 
Its inherent smoothness ensures that gradients are well behaved, 
facilitating stable optimization during training, 
while local control allows for fine grained adjustments in specific 
regions of the input space without affecting the entire function. 
In addition, geometric interpretability provides insight into 
how the function behaves across different regions and its capacity to 
represent both polynomial and non-polynomial shapes with a relatively 
small number of control points makes it highly efficient. 
Due to its theoretical universality, \texttt{NELE} can, in principle, 
approximate any continuous function to arbitrary accuracy, 
making it a flexible candidate for general-purpose activation in neural networks. 
To evaluate its practical performance, \texttt{NELE} was tested across multiple 
state of the art benchmark datasets spanning diverse machine 
learning tasks. The empirical results consistently demonstrate that 
\texttt{NELE} achieves superior performance relative to existing activation 
function methods, both in terms of convergence speed and predictive accuracy, 
highlighting its potential as a universal activation function capable of 
capturing complex functional relationships.
%%%%%%%%%%%%%%%%%%%%%%%%%%%%%%%%%%%%%%%%%%%%%%%%%%%%%%%%%%%%%%%%%%%%%%%%%%
%%
%% The acknowledgments section is defined using the "acks" environment
%% (and NOT an unnumbered section). This ensures the proper
%% identification of the section in the article metadata, and the
%% consistent spelling of the heading.
\begin{acks}
This work was carried out independently by the author.
\end{acks}

%%
%% The next line prints the references.
\printbibliography

%%
%% If your work has an appendix, this is the place to put it.
\appendix
\section{Function analysis}
\begin{figure}[H]\label{fig:afcollect}
  \begin{center}
    \includegraphics[width=0.6\linewidth]{../../analysis/analysis.pdf}
  \caption{The solid lines represent the functions $y(x)$ for different activation functions. 
The points of maximum curvature, $\kappa_{\mathrm{max}}$, are indicated by markers along each curve. 
The dashed lines represent the curvature, $\kappa(x)$, as a function of $x$. 
Each color corresponds to a different activation function as shown in the legend.}
  \end{center}
  \Description{Analysis}
\end{figure}
\subsection{Tanh}\label{app:tanh}

\begin{displaymath}
  y = \tanh{x}
\end{displaymath}
\begin{displaymath}
y' = \operatorname{sech}^2(x) = 1 - y^2
\end{displaymath}
\begin{displaymath}
y'' = -2\operatorname{sech}^2(x)\tanh{x} = -2y(1 - y^2)
\end{displaymath}
\begin{displaymath}
\kappa = \frac{|-2\operatorname{sech}^2(x)\tanh{x}|}{(1 + (\operatorname{sech}^2(x))^2)^\frac{3}{2}}
\end{displaymath}
\begin{displaymath}
\frac{d \kappa}{dx} = 0 
\end{displaymath}

\begin{displaymath}
\text{Using chain rule}, \frac{d \kappa}{dx} = \frac{d \kappa}{dy}\frac{dy}{dx}\text{. Then } \frac{d \kappa}{dy} = 0 \text{ or } 1 - y^2 = 0
\end{displaymath}
\begin{displaymath}
  \frac{d}{dy}\left( \frac{2y(1 - y^2)}{(1 + (1 - y^2)^2)^{3/2}}\right) = 3y^6 - 5y^4 - 2y^2 + 2 = 0
\end{displaymath}
\begin{displaymath}
  \text{Solving: }y = 0.9196, \kappa_{\max} = 0.5073
\end{displaymath}

\subsection{Sigmoid}\label{app:sigmoid}

\begin{displaymath}
  y = \frac{1}{1 + e^{-x}}
\end{displaymath}
\begin{displaymath}
  y' = \frac{-e^{-x}}{(1 + e^{-x})^2}
\end{displaymath}
\begin{displaymath}
  y'' = \frac{e^{-x}(1 - e^{-x})}{(1 + e^{-x})^3}
\end{displaymath}
\begin{displaymath}
  \kappa = \frac{\frac{e^{-x}(1 - e^{-x})}{(1 + e^{-x})^3}}{(1 + ( \frac{-e^{-x}}{(1 + e^{-x})^2})^2)^\frac{3}{2}}
\end{displaymath}
Solving this equation leads to a transcendental condition, 
which has no simple closed-form solution, so the 
locations of maximum curvature must be obtained numerically.

\subsection{Softplus}\label{app:softplus}

\begin{displaymath}
  y = \ln{({1 + e^{x}})}
\end{displaymath}
\begin{displaymath}
  y' = \frac{e^x}{1 + e^x}
\end{displaymath}
\begin{displaymath}
  y'' = \frac{e^{x}}{(1 + e^{x})^2}
\end{displaymath}
\begin{displaymath}
  \kappa = \frac{\frac{e^{x}}{(1 + e^{x})^2}}{(1 + (\frac{e^x}{1 + e^x})^2)^\frac{3}{2}}
\end{displaymath}
\begin{displaymath}
  \text{Let }s = \frac{e^x}{1 + e^x}, \kappa = \frac{s(1 - s)}{(1 + s^2)^\frac{3}{2}}
\end{displaymath}
\begin{displaymath}
  \frac{d\kappa}{dx} = 0, \frac{d}{dx}\bigg(\frac{s(1 - s)}{(1 + s^2)^\frac{3}{2}}\bigg) = 0
\end{displaymath}
\begin{displaymath}
  s = -1, 0.38197, 2.6180 \text{. Substituting, }x = -0.5, \kappa_{\text{max}} = 0.1913
\end{displaymath}

\subsection{ELU}\label{app:elu}
 \begin{displaymath}
  y = \begin{cases}
  x & \text{if } x \geq 0, \\
  \alpha (e^x - 1) & \text{if } x < 0
\end{cases}
 \end{displaymath}

 \begin{displaymath}
  y' = \begin{cases}
  1 & \text{if } x \geq 0, \\
  \alpha e^x & \text{if } x < 0
\end{cases}
 \end{displaymath}
\begin{displaymath}
  y'' = \begin{cases}
  0 & \text{if } x \geq 0, \\
  \alpha e^x & \text{if } x < 0
\end{cases}
\end{displaymath}
\begin{displaymath}
  \kappa = \begin{cases}
  0 & \text{if } x \geq 0, \\
  \frac{\alpha e^x}{(1 + \alpha^2 e^{2x})^\frac{3}{2}}  & \text{if } x < 0
\end{cases}
\end{displaymath}
\begin{displaymath}
  \frac{d \kappa}{dx} = 0
\end{displaymath}
\begin{displaymath}
  \frac{d}{dx} \left( \frac{\alpha e^x}{(1 + \alpha^2 e^{2x})^{3/2}} \right)
= \frac{\alpha e^x \left( 1 - 2 \alpha^2 e^{2x} \right)}{(1 + \alpha^2 e^{2x})^{5/2}} = 0
\end{displaymath}
\begin{displaymath}
  \alpha e^x \left( 1 - 2 \alpha^2 e^{2x} \right) = 0
\end{displaymath}
\begin{displaymath}
  e^x = 0 \text{ has no solution} \text{, therefore, } x = -\frac{1}{2} \ln(2 \alpha^2)
\end{displaymath}
\begin{displaymath}
  \text{Substituting back, }\kappa_{\text{max}} = \frac{2}{3\sqrt3}
\end{displaymath}

\subsection{SiLU}\label{app:silu}

\begin{displaymath}
y = \frac{x}{1 + e^{-x}}
\end{displaymath}
\begin{displaymath}
y' = \frac{1 + e^{-x} + xe^{-x}}{(1 + e^{-x})^2}
\end{displaymath}
\begin{displaymath}
y'' = \frac{e^{-x}(2 - x) + e^{-2x}(x+2)}{(1 + e^{-x})^3}
\end{displaymath}
\begin{displaymath}
\kappa = \frac{\frac{e^{-x}(2 - x) + e^{-2x}(x+2)}{(1 + e^{-x})^3}}{(1 + (\frac{1 + e^{-x} + xe^{-x}}{(1 + e^{-x})^2})^2)^\frac{3}{2}}
= \frac{(1 + e^{-x})^2(e^{-x}(2 - x) + e^{-2x}(x+2))}{[(1 + e^{-x})^4 + (1 + e^{-x} + xe^{-x})^2]^\frac{3}{2}}
\end{displaymath}
\begin{displaymath}
\frac{d \kappa}{dx} = 0
\end{displaymath}
\begin{displaymath}
\frac{d}{dx}\biggr(\frac{(1 + e^{-x})^2(e^{-x}(2 - x) + e^{-2x}(x+2))}{[(1 + e^{-x})^4 + (1 + e^{-x} + xe^{-x})^2]^\frac{3}{2}}\biggr) = 0
\end{displaymath}
Using quotient rule and consider the first term in the numerator after differentiation :
$[e^{-x}(x-3)-e^{-2x}(2x+3)](1+e^{-x})^2$. This leads to a nonlinear 
transcendental equation. Since, they do not simplify into algebraic 
forms that allow symbolic isolation of x, an exact analytical solution for the location of 
maximum curvature is not possible and numerical or graphical methods must be used.

\subsection{GELU}\label{app:gelu}

\begin{displaymath}
y = x \phi(x) \approx 0.5x(1 + \tanh [\sqrt{2 / \pi}(x + 0.44715 x^3)])
\end{displaymath}
\begin{displaymath}
\text{Let }k = \sqrt{\frac{2}{\pi}}, a = 0.441715 \text{ and }u(x) = k(x + ax^3)
\end{displaymath}
\begin{displaymath}
y' = 0.5(1 + \tanh{u}) + 0.5ax\operatorname{sech}^2(u)(1 + 3ax^2)
\end{displaymath}
\begin{displaymath}
y'' = 0.5a(1 + 3ax^2)\operatorname{sech}^2(u) + 0.5x[-2\operatorname{sech}^2(u)\tanh{u}(a(1 + 3ax^2))^2 + 6kax\operatorname{sech}^2(u)]
\end{displaymath}

\subsection{Mish}\label{app:mish}
\begin{displaymath}
y = x \cdot \tanh\!\left(\ln(1 + e^{x})\right)
\end{displaymath}
\begin{displaymath}
y' = \tanh{(\ln{(1 + e^x)})} + \frac{xe^x}{1 + e^x}\operatorname{sech}^2(\ln{(1 + e^x)})
\end{displaymath}
\begin{displaymath}
  y'' = a + b + c + d
\end{displaymath}
\begin{displaymath}
  \text{where, } a = \frac{e^x}{1 + e^x}\operatorname{sech}^2(\ln{(1 + e^x)})
\end{displaymath}
\begin{displaymath}
  b = \frac{\operatorname{sech}^2(\ln{(1 + e^x)})}{1+e^{-x}} 
\end{displaymath}
\begin{displaymath}
  c = \frac{xe^{-x}}{(1 + e^{-x})^2} \operatorname{sech}^2(\ln{(1 + e^x)})
\end{displaymath}
\begin{displaymath}
  d = \frac{-2x}{(1 + e^{-x})^2} \operatorname{sech}^2(\ln{(1 + e^x)} \tanh{\ln{(1 + e^x)}}
\end{displaymath}

\subsection{LeLU}\label{app:LeLU}

\begin{displaymath}
y = e^{(1 - \beta)x} + 1 - \beta x
\end{displaymath}
\begin{displaymath}
y' = (1 - \beta)e^{(1 - \beta)x} - \beta
\end{displaymath}
\begin{displaymath}
y'' = (1 - \beta)^2e^{(1 - \beta)x}
\text{Let }1 - \beta = a, e^{(1 - \beta)x} = e^{ax} = b
\end{displaymath}
\begin{displaymath}
\kappa = \frac{a^2b}{(1 + a^2b^{2})^\frac{3}{2}}
\end{displaymath}
\begin{displaymath}
\frac{d \kappa}{dx} = 0
\end{displaymath}
\begin{displaymath}
\text{Solving, }\beta = 1 \pm \sqrt{\frac{1}{2}}, 1, \text{Substituting }x = 0, \pm \frac{\sqrt{2}}{2}\ln2
\end{displaymath}

\section{Curve approximation with \texttt{NELE}\label{app:nurbsparams}}
To approximate a target curve with \texttt{NELE}, an optimization problem is formulated to minimize the mean squared error (MSE) between the target function and the NURBS representation.
Optimization algorithm L-BFGS-B \cite{scipy_optimize_minimize} is used to efficiently navigate the parameter space, 
producing a smooth curve that closely follows the desired shape.
\begin{table}[H]
  \caption{Control point coordinates and their respective weights for approximation on the negative space of AF using \texttt{NELE}}
  \label{tab:curveapprox}
  \begin{tabular}{@{}llllll@{}}
    \toprule
    AF&$p_0$&$p_1$&$p_2$&$p_3$&$w$\\
    \midrule
    \texttt{Tanh} & [-4.0, -1.0] &    [-0.884, -0.983] &    [-1.099, -1.078] &    [0.0, -0.001] &    [1.942 , 0.831 , 0.253 , 0.209]\\
    \texttt{Sigmoid} & [-4.0, 0.017] &    [-2.03, 0.059] &    [-1.498, 0.112] &    [0.0, 0.501] &    [1.31 , 0.886 , 0.719 , 0.808]\\
    \texttt{Softplus} & [-4.0, 0.017] &    [-2.332, 0.056] &    [-1.157, 0.105] &    [-0.0, 0.694] &    [1.233 , 0.985 , 0.808 , 0.701]\\
    \texttt{ELU} & [-4.0, -0.984] &    [-2.394, -0.935] &    [-0.913, -0.912] &    [0.001, -0.001] &    [1.227 , 1.017 , 0.79 , 0.54]\\
    \texttt{SiLU} & [-4.0, -0.078] &    [-2.383, -0.112] &    [-1.011, -0.547] &    [-0.0, 0.003] &    [1.226 , 1.012 , 0.806 , 0.611]\\
    \texttt{LeLU} & [-4.0, -1.373] &    [-2.375, -1.166] &    [-0.898, -0.917] &    [-0.001, 0.002] &    [1.235 , 1.015 , 0.776 , 0.524]\\
    \texttt{Mish} & [-4.0, -0.079] &    [-2.358, -0.113] &    [-0.868, -0.598] &    [-0.001, 0.008] &    [1.275 , 1.037 , 0.779 , 0.509]\\
    \texttt{GELU} & [-4.0, 0.0] &    [-0.321, 0.058] &    [-2.134, -0.913] &    [0.0, 0.0] &    [1.980 , 0.717 , 0.124 , 0.198]\\
  \bottomrule
\end{tabular}
\end{table}

\section{Online Resources}

The source code used in this paper are publicly available to ensure reproducibility and facilitate further research:

\begin{itemize}
    \item \textbf{\texttt{NELE} Implementation:} Python implementation of the proposed NELE activation function, including training scripts and examples. Available at: 
    \url{https://github.com/DjentleViBe/NELE}.
    
    \item \textbf{Additional Figures:} Addiitonal plots for the curve approximation from Table ~\ref{tab:curveapprox} are provided in the repository.

    \item \textbf{Settings:} Parameters used for reproducing the results in this paper is provided in the
    \texttt{README} file of the repository.
  
\end{itemize}

It is recommended that users follow the provided \texttt{README} instructions for replicating experiments. This repository will remain publicly accessible for the foreseeable future.

\section{Reproducibility Checklist for JAIR}

Select the answers that apply to your research -- one per item. 

\subsection*{All articles:}

%\hh{revised for stylistic consistency:}
\begin{enumerate}
    \item All claims investigated in this work are clearly stated. 
    [yes]
    \item Clear explanations are given how the work reported substantiates the claims. 
    [yes]
    \item Limitations or technical assumptions are stated clearly and explicitly. 
    [yes]
    \item Conceptual outlines and/or pseudo-code descriptions of the AI methods introduced in this work are provided, and important implementation details are discussed. 
    [yes]
    \item 
    Motivation is provided for all design choices, including algorithms, implementation choices, parameters, data sets and experimental protocols beyond metrics.
    [yes/partially/no]
\end{enumerate}

\subsection*{Articles containing theoretical contributions:}
Does this paper make theoretical contributions? 
[yes/no] 

If yes, please complete the list below.

\begin{enumerate}
    \item All assumptions and restrictions are stated clearly and formally. 
    [yes/partially/no]
    \item All novel claims are stated formally (e.g., in theorem statements). 
    [yes/partially/no]
    \item Proofs of all non-trivial claims are provided in sufficient detail to permit verification by readers with a reasonable degree of expertise (e.g., that expected from a PhD candidate in the same area of AI). [yes/partially/no]
    \item
    Complex formalism, such as definitions or proofs, is motivated and explained clearly.
%hh: was:
%Proof sketches or intuitions are given for complex and/or novel results.
    [yes/partially/no]
    \item 
    The use of mathematical notation and formalism serves the purpose of enhancing clarity and precision; gratuitous use of mathematical formalism (i.e., use that does not enhance clarity or precision) is avoided.
    [yes/partially/no]
    \item 
    Appropriate citations are given for all non-trivial theoretical tools and techniques. 
    [yes/partially/no]
\end{enumerate}

\subsection*{Articles reporting on computational experiments:}
Does this paper include computational experiments? [yes/no]

If yes, please complete the list below.
\begin{enumerate}
    \item 
    All source code required for conducting experiments is included in an online appendix 
    or will be made publicly available upon publication of the paper.
    The online appendix follows best practices for source code readability and documentation as well as for long-term accessibility.
    [yes/partially/no]
    \item The source code comes with a license that
    allows free usage for reproducibility purposes.
    [yes/partially/no]
    \item The source code comes with a license that
    allows free usage for research purposes in general.
    [yes/partially/no]
    \item 
    Raw, unaggregated data from all experiments is included in an online appendix 
    or will be made publicly available upon publication of the paper.
    The online appendix follows best practices for long-term accessibility.
    [yes/partially/no]
    \item The unaggregated data comes with a license that
    allows free usage for reproducibility purposes.
    [yes/partially/no]
    \item The unaggregated data comes with a license that
    allows free usage for research purposes in general.
    [yes/partially/no]
    \item If an algorithm depends on randomness, then the method used for generating random numbers and for setting seeds is described in a way sufficient to allow replication of results. 
    [yes/partially/no/NA]
    \item The execution environment for experiments, the computing infrastructure (hardware and software) used for running them, is described, including GPU/CPU makes and models; amount of memory (cache and RAM); make and version of operating system; names and versions of relevant software libraries and frameworks. 
    [yes/partially/no]
    \item 
    The evaluation metrics used in experiments are clearly explained and their choice is explicitly motivated. 
    [yes/partially/no]
    \item 
    The number of algorithm runs used to compute each result is reported. 
    [yes/no]
    \item 
    Reported results have not been ``cherry-picked'' by silently ignoring unsuccessful or unsatisfactory experiments. 
    [yes/partially/no]
    \item 
    Analysis of results goes beyond single-dimensional summaries of performance (e.g., average, median) to include measures of variation, confidence, or other distributional information. 
    [yes/no]
    \item 
    All (hyper-) parameter settings for 
    the algorithms/methods used in experiments have been reported, along with the rationale or method for determining them. 
    [yes/partially/no/NA]
    \item 
    The number and range of (hyper-) parameter settings explored prior to conducting final experiments have been indicated, along with the effort spent on (hyper-) parameter optimisation. 
    [yes/partially/no/NA]
    \item 
    Appropriately chosen statistical hypothesis tests are used to establish statistical significance
    in the presence of noise effects.
    [yes/partially/no/NA]
\end{enumerate}


\subsection*{Articles using data sets:}
Does this work rely on one or more data sets (possibly obtained from a benchmark generator or similar software artifact)? 
[yes/no]

If yes, please complete the list below.
\begin{enumerate}
    \item 
    All newly introduced data sets 
    are included in an online appendix 
    or will be made publicly available upon publication of the paper.
    The online appendix follows best practices for long-term accessibility with a license
    that allows free usage for research purposes.
    [yes/partially/no/NA]
    \item The newly introduced data set comes with a license that
    allows free usage for reproducibility purposes.
    [yes/partially/no]
    \item The newly introduced data set comes with a license that
    allows free usage for research purposes in general.
    [yes/partially/no]
    \item All data sets drawn from the literature or other public sources (potentially including authors' own previously published work) are accompanied by appropriate citations.
    [yes/no/NA]
    \item All data sets drawn from the existing literature (potentially including authors’ own previously published work) are publicly available. [yes/partially/no/NA]
    %\item All data sets that are not publicly available are described in detail.
    %[yes/partially/no/NA]
    \item All new data sets and data sets that are not publicly available are described in detail, including relevant statistics, the data collection process and annotation process if relevant.
    [yes/partially/no/NA]
    \item 
    All methods used for preprocessing, augmenting, batching or splitting data sets (e.g., in the context of hold-out or cross-validation)
    are described in detail. [yes/partially/no/NA]
\end{enumerate}

\subsection*{Explanations on any of the answers above (optional):}

[Text here; please keep this brief.]

\end{document}
\endinput
%%
%% End of file `sample-acmlarge.tex'.
